\documentclass[11pt, a4paper]{article}

\usepackage[T1]{fontenc}
\usepackage[utf8]{inputenc}
\usepackage[a4paper, top=2.5cm, bottom=2.5cm, left=2.5cm, right=2.5cm]{geometry}
\usepackage{booktabs}
\usepackage{enumitem}
\usepackage[hidelinks]{hyperref}
\usepackage{parskip}
\usepackage{amsmath}
\usepackage{listings}
\usepackage{xcolor}

\definecolor{codigobg}{RGB}{245,245,245}
\definecolor{keyword}{RGB}{0,100,150}
\lstset{language=R, basicstyle=\ttfamily\small, backgroundcolor=\color{codigobg},
  keywordstyle=\color{keyword}\bfseries, frame=single, breaklines=true, showstringspaces=false}

\setlength{\parindent}{0pt}

\pagestyle{plain}

\begin{document}

\begin{center}
{\large\bfseries PONTIFICIA UNIVERSIDAD JAVERIANA}\\[3pt]
Facultad de Ciencias Econ\'omicas y Administrativas\\[2pt]
Estad\'istica 2026-II --- Prof.\ Jaime Polanco Jim\'enez\\[6pt]
{\LARGE\bfseries Problem Set 1}\\[3pt]
{\large Estad\'istica Descriptiva}\\[6pt]
Asignado: Julio 31 (Sesi\'on 1) \quad$|$\quad Entrega: Agosto 14 (Sesi\'on 2)
\end{center}

\vspace{6pt}
\rule{\linewidth}{0.5pt}
\vspace{2pt}
{\centering\small Repositorio: \url{https://github.com/polanco-jaime/Curso-Estadistica-2026-3}\par}
\vspace{6pt}

\textbf{Instrucciones.} Este Problem Set es individual y se resuelve en R.
\textbf{Fecha l\'imite de entrega: el d\'ia anterior a la sesi\'on indicada, a m\'as tardar a las 11:59\,p.m.\ (hora Colombia, UTC$-$5).} Entregas tard\'ias pierden el bono. El c\'odigo R debe incluirse.

\section*{Datos}

Los datos se encuentran en el repositorio del curso en formato Parquet:

\begin{lstlisting}
library(arrow)
saber_pro <- read_parquet("datos/saber_pro/saber_pro.parquet")
\end{lstlisting}

\textbf{Diccionario de variables utilizadas:}

\vspace{4pt}
\begin{tabular}{l l l}
\toprule
\textbf{Variable} & \textbf{Descripci\'on} & \textbf{Tipo} \\
\midrule
\texttt{MOD\_INGLES\_PUNT}       & Puntaje m\'odulo de ingl\'es (Saber Pro) & Num\'erico \\
\texttt{ESTU\_DEPTO\_RESIDE}     & Departamento de residencia del estudiante & Categ\'orico \\
\texttt{INST\_ORIGEN}            & Tipo de IES: Oficial / Privada & Categ\'orico \\
\texttt{FAMI\_ESTRATOVIVIENDA}   & Estrato socioecon\'omico (1--6) & Ordinal \\
\bottomrule
\end{tabular}

\section*{Enunciado}

Realice un an\'alisis descriptivo completo de la variable \texttt{MOD\_INGLES\_PUNT}
(puntaje de ingl\'es en Saber Pro).
Su an\'alisis debe incluir:

\begin{enumerate}[label=\alph*)]
  \item \textbf{Medidas de tendencia central:} calcule la media, mediana y moda de
        \texttt{MOD\_INGLES\_PUNT} para el conjunto completo de datos.

  \item \textbf{Medidas de dispersi\'on:} calcule la desviaci\'on est\'andar, varianza,
        rango y rango intercuart\'ilico (IQR).

  \item \textbf{Forma de la distribuci\'on:} calcule e interprete los coeficientes de
        asimetr\'ia y curtosis. ¿La distribuci\'on es sim\'etrica? ¿Es leptoc\'urtica
        o platic\'urtica?

  \item \textbf{Cuartiles:} presente la tabla de cuartiles (Q1, Q2, Q3) y percentiles
        relevantes (P10, P90).

  \item \textbf{Desagregaci\'on:} repita los c\'alculos anteriores desagregando por:
        \begin{itemize}
          \item Departamento (\texttt{ESTU\_DEPTO\_RESIDE})
          \item Tipo de IES: p\'ublica vs.\ privada (\texttt{INST\_ORIGEN})
          \item Estrato socioecon\'omico (\texttt{FAMI\_ESTRATOVIVIENDA})
        \end{itemize}

  \item \textbf{Visualizaciones:} produzca un m\'inimo de 5 gr\'aficos en R
        (\texttt{ggplot2}) que complementen su an\'alisis. Sugerencias:
        \begin{itemize}
          \item Histograma de \texttt{MOD\_INGLES\_PUNT}
          \item Box plot por tipo de IES
          \item Box plot por estrato
          \item Gr\'afico de barras de media por departamento
          \item Densidad (KDE) comparando p\'ublica vs.\ privada
        \end{itemize}

  \item \textbf{Conclusiones:} en un p\'arrafo breve, resuma los hallazgos
        principales de su an\'alisis descriptivo.
\end{enumerate}

\vspace{10pt}
\rule{\linewidth}{0.4pt}
\begin{center}
{\small Cada entrega completa suma $+0.0625$ a la nota final.}
\end{center}

\end{document}
